\chapter{ПРОЕКТИРОВАНИЕ И РЕАЛИЗАЦИЯ ВЫЧИСЛИТЕЛЬНОГО ЯДРА}

\section{Общая архитектура программного комплекса}

Разрабатываемый программный комплекс реализован по архитектурному принципу «одно вычислительное ядро и несколько каналов доставки результата».  
Вычисления термодинамических свойств сосредоточены в единственном компоненте \texttt{if97\_core}.  
Пользовательские приложения и сервисный контур не дублируют формулы, а вызывают единый API ядра и отображают или передают результат.  

Такое разделение позволяет одновременно решать две задачи.  
Первая задача --- обеспечить одинаковые результаты вычислений во всех режимах использования (desktop, web и сервис).  
Вторая задача --- локализовать математическую сложность стандарта IF97 в одном месте и повысить управляемость изменений.  

Для обмена данными между подсистемами используется общий слой DTO \texttt{if97\_app\_api}.  
Он фиксирует формат входных параметров и итоговых свойств, снижая связность между ядром, интерфейсами и сервисом.  
Реализация выполнена на языке Rust, что обеспечивает строгую типизацию, выразительные модели ошибок и возможность получения высокой производительности без потери надёжности \cite{rust_book}.  

\section{Структура и модули \texttt{if97\_core}}

Внутренняя организация \texttt{if97\_core} строится как набор слоёв, каждый из которых решает ограниченный класс задач.  
Подход к декомпозиции соответствует практикам DDD, где выделяются доменные типы и доменные правила, а вычисления оформляются как детерминированные stateless-операции.  

Типовая структура ядра включает следующие функциональные блоки:
\begin{itemize}
\item \texttt{domain} --- типы предметной области, единицы измерения, ошибки, итоговое состояние воды.  
\item \texttt{engine} --- валидация входа, классификация области/региона и запуск «плана вычисления».  
\item \texttt{regions} --- реализации уравнений IF97 для регионов 1--5, а также обратные зависимости и численные процедуры там, где это необходимо.  
\item \texttt{topology} --- фазовые границы, линия насыщения (Region~4), граница B23 и правила выбора региона.  
\item \texttt{calculator} --- публичный фасад \texttt{If97}, который принимает типизированный вход и возвращает унифицированный результат.  
\end{itemize}

Отдельно выделяется подсистема констант и таблиц коэффициентов IF97, которая подключается к региональным моделям как неизменяемые данные.  
Для сценариев визуализации используется специализированный API насыщения, позволяющий строить границы купола без лишних вычислений двухфазной смеси.  

\section{Типобезопасные единицы измерения и модель данных}

В ядре используется типобезопасное представление физических величин.  
Давление, температура, энтальпия и другие параметры задаются через newtype-обёртки, которые явно фиксируют размерность.  
Например, давление передаётся как \texttt{MegaPascal}, а температура как \texttt{Kelvin}.  

Такой подход решает несколько инженерных задач.  
Во-первых, он предотвращает ошибки смешения единиц измерения и уменьшает вероятность некорректного вызова API.  
Во-вторых, он делает интерфейс читаемым и самодокументируемым, поскольку размерность видна уже в сигнатуре функций.  
В-третьих, он упрощает внедрение централизованной валидации диапазонов, поскольку проверка привязана к конкретному доменному типу.  

Результатом вычисления выступает агрегированное состояние воды \texttt{WaterState}.  
Оно включает входные параметры и основные расчётные свойства, востребованные в инженерных задачах:
\begin{itemize}
\item давление $p$ и температура $T$;  
\item удельный объём $v$ и плотность $\rho$;  
\item энтальпия $h$, энтропия $s$ и внутренняя энергия $u$;  
\item изобарная теплоёмкость $c_p$ и скорость звука $w$;  
\item номер региона IF97, в котором выполнялся расчёт.  
\end{itemize}

Унифицированная структура результата обеспечивает одинаковое представление состояния для всех маршрутов расчёта и всех внешних потребителей ядра.  

\section{Механизм маршрутизации расчётов и классификации регионов}

Стандарт IAPWS-IF97 использует региональные уравнения, поэтому основой вычислительного процесса является корректное определение региона и выбор соответствующей модели \cite{iapws_if97_2012}.  
В ядре эта логика реализована в виде маршрутизатора, который анализирует входные параметры и выбирает план вычисления.  

Публичный фасад \texttt{If97} поддерживает несколько основных маршрутов:
\begin{itemize}
\item $pT$ --- прямой расчёт по давлению и температуре;  
\item $pH$ и $pS$ --- обратные расчёты по давлению и энтальпии или по давлению и энтропии;  
\item $pX$ --- расчёт в двухфазной области по давлению и степени сухости $x$;  
\item $\rho T$ --- расчёт по плотности и температуре как численное расширение, в котором подбирается давление, после чего используется стандартный маршрут $pT$.  
\item метастабильный пар --- специализированный режим, ограниченный областью применимости (расширение Region~2).  
\end{itemize}

Внутри каждого маршрута вычисления сводятся к единому сценарию.  
Сначала выполняется валидация диапазонов и проверка физической допустимости входа.  
Далее определяется регион IF97 с учётом топологии области, линии насыщения (Region~4) и границы B23.  
После этого выбирается алгоритм вычислений: прямые формулы, обратные аналитические зависимости или численный решатель.  

Для обратных маршрутов по парам $(p,h)$ и $(p,s)$ используются обратные формулы IF97 там, где они даны стандартом, что позволяет избежать итераций.  
Для задач, связанных с Region~3, используются как базовая формулировка через потенциал Гельмгольца, так и дополнительные обратные зависимости, опубликованные IAPWS (например, $v(p,T)$ для Region~3) \cite{iapws_vpt_region3}.  
В околокритической области Region~3 при реализации обратных маршрутов возникают численные трудности, поскольку прямое решение систем уравнений становится неустойчивым при сингулярности Якобиана.  
Поэтому для стабилизации расчётов применяется 2D-решатель Левенберга--Марквардта с предварительным сеточным поиском и демпфированием шага.  

\section{Генерация таблиц коэффициентов и организация констант}

В IAPWS-IF97 большая часть формул представляет собой ряды, параметризованные наборами коэффициентов.  
Чтобы исключить ошибки ручного ввода и не выполнять файловый ввод-вывод при запуске, коэффициенты хранятся в CSV-таблицах и преобразуются в статический код на этапе сборки.  

На практике используется скрипт \texttt{build.rs}, который читает CSV, нормализует порядок данных и генерирует Rust-модуль со статическими массивами коэффициентов.  
Сгенерированный код подключается через \texttt{include!} и становится частью итогового бинарного файла.  

Данный подход обеспечивает следующие преимущества:
\begin{itemize}
\item отсутствие runtime I/O при обращении к коэффициентам;  
\item воспроизводимость сборки и предсказуемость поведения при развертывании;  
\item ускорение вычислений за счёт обращения к константам вместо динамических структур;  
\item снижение риска ошибок из-за несовместимых форматов или различий в окружении.  
\end{itemize}

\section{Обработка ошибок и валидация}

Для инженерного инструмента важно не только вычислить свойства, но и корректно обработать случаи выхода за область применимости стандарта и некорректные комбинации входных параметров.  
Поэтому в \texttt{if97\_core} выделена централизованная валидация входа и единая модель ошибок предметной области.  

На этапе валидации проверяются:
\begin{itemize}
\item принадлежность давления и температуры допустимым диапазонам IF97;  
\item положительность плотности;  
\item диапазон степени сухости $x \in [0,1]$;  
\item специальные ограничения для метастабильного пара и обратных маршрутов.  
\end{itemize}

Ошибки рассматриваются как часть предметной логики и возвращаются вызывающей стороне в диагностируемом виде.  
Паники и неявные «тихие» состояния (например, распространение \texttt{NaN} в результат) в корректных сценариях исключаются за счёт строгих проверок и типизации.  
Поскольку ядро используется всеми каналами доставки, единая модель ошибок обеспечивает одинаковое поведение в desktop-приложении, web-оболочке и сервисе.  

\section{Верификация ядра и тестовые сценарии}

Верификация ядра включает модульные и интеграционные тесты, а также сравнение результатов с контрольными значениями.  
Контрольные точки берутся из материалов стандарта IAPWS-IF97 и проверяют корректность как прямых, так и обратных расчётов \cite{iapws_if97_2012}.  

Для автоматической проверки используются:
\begin{itemize}
\item unit-тесты региональных формул и вспомогательных алгоритмов;  
\item интеграционные тесты маршрутов $pT$, $pH$, $pS$, $pX$ и $\rho T$;  
\item сценарии сверки по CSV-таблицам контрольных значений (регрессионная проверка).  
\end{itemize}

Помимо корректности, для практической применимости важна производительность.  
Поэтому отдельными сценариями оценивается скорость пакетных вычислений и стоимость построения данных для диаграмм.  
По результатам нагрузочных проверок пропускная способность пакетного расчёта находится на уровне сотен тысяч строк в секунду в release-сборке, а вычисление купола насыщения укладывается в десятки миллисекунд на диаграмму.  
